\documentclass[a4paper,10pt]{article}

% ---------- Packages ----------
\usepackage[margin=0.45in]{geometry}
\usepackage{titlesec}
\usepackage{enumitem}
\usepackage{hyperref}
\usepackage{tgtermes}
\usepackage[T1]{fontenc}
\usepackage{xcolor}
\usepackage{tabularx}
\usepackage{array}

% ---------- Page Setup ----------
\pagestyle{empty}
\setlength{\parindent}{0pt}
\setlength{\parskip}{0pt}

\hypersetup{
    colorlinks=true,
    urlcolor=blue,
    linkcolor=blue,
    citecolor=blue
}

% ---------- Section Formatting ----------
\titleformat{\section}
  {\Large\scshape}
  {}
  {0em}
  {}
  [\titlerule]

\titlespacing*{\section}{0pt}{3pt}{3pt}

% ---------- List Formatting ----------
\setlist[itemize]{
    leftmargin=1.2em,
    itemsep=1pt,
    topsep=2pt,
    parsep=0pt,
    partopsep=0pt
}

\renewcommand\labelitemi{$\bullet$}

% ---------- Custom Commands ----------
\newcommand{\resumeHeader}[8]{
    \begin{center}
        {\LARGE \textbf{#1}}\\
        \href{mailto:#2}{#2} \,|\, #3 \,|\, \href{#4}{LinkedIn} \,|\, \href{#5}{GitHub} \,|\, \href{#6}{LeetCode}\\
        #7\\
        #8
    \end{center}
}

\newcommand{\resumeSubheading}[4]{
    \textbf{#1} \hfill #2\\
    \textit{#3} \hfill #4
}

\newcommand{\resumeProjectHeading}[3]{
    \textbf{#1}\\
    \textit{#2} \hfill #3
}

\newcommand{\resumeEducationHeading}[4]{
    \textbf{#1} \hfill #2\\
    #3\\
    #4
}

\newcommand{\skillLine}[2]{%
    \noindent\textbf{#1:} #2\par\vspace{0pt}%
}

% ---------- Document ----------
\begin{document}
\vspace*{-35pt}

% ---------- Header ----------
\resumeHeader
{Prerak Patel}
{patel.prerak2798@gmail.com}
{+1 (321) 610-0582}
{https://www.linkedin.com/in/prerakpatel51020021998/}
{https://github.com/prerakpatel51}
{https://leetcode.com/u/PrerakPatel510/}
{Machine Learning Engineer \,|\, Deep Learning \,|\, Python}
{Melbourne, Florida \,|\, Permanent Resident, U.S. Person}

\vspace{-8pt}

% ---------- Summary ----------
\section*{Summary}
Machine Learning Engineer with hands-on experience building end-to-end ML workflows from data ingestion and preprocessing to training, evaluation, and scalable execution. Strong in \textbf{Python}, \textbf{PyTorch}, time-series forecasting, and generative modeling (VAE, Diffusion, Conditional Flow Matching) on large geospatial and radar/satellite datasets. Delivered reproducible pipelines with experiment tracking, multi-GPU training on Slurm HPC, and latency/performance improvements in backend systems. Comfortable working across research-to-product engineering with \textbf{AWS}, \textbf{Docker}, and modern ML tooling.

% ---------- Skills ----------
\section*{Skills}
\skillLine{Languages}{\textbf{Python}, C++, Java, SQL, JavaScript}
\skillLine{Machine Learning}{\textbf{Machine Learning}, \textbf{Deep Learning}, Supervised Learning, Time-Series Forecasting, NLP, Sentiment Analysis, Feature Engineering, Model Evaluation}
\skillLine{Frameworks}{\textbf{PyTorch}, TensorFlow, Keras, Scikit-Learn}
\skillLine{Data \& Visualization}{NumPy, Pandas, Polars, Dask, SciPy, Statsmodels, Matplotlib, Plotly, OpenCV}
\skillLine{MLOps / DevOps}{\textbf{AWS}, \textbf{Docker}, Kubernetes, Git, CI/CD, MLflow, Hugging Face Spaces, Gradio}
\skillLine{Backend}{Django, RESTful APIs, Redis, PostgreSQL}
\skillLine{Compute}{Slurm, GPU Clusters, Multi-GPU Training}
\vspace{-2pt}

% ---------- Work Experience ----------
\section*{Work Experience}

\resumeSubheading
{Florida Institute of Technology, Melbourne, FL}
{April 2026 -- Present}
{Research Assistant}
{Advisor: Dr. Yang}
\begin{itemize}
    \item Building \textbf{Python} ML and time-series pipelines for cryptocurrency and stock market forecasting using NLP preprocessing, feature engineering, regression/classification, and model evaluation.
    \item Performing financial sentiment analysis on news/social/financial text to study relationships between sentiment signals and asset price movement.
\end{itemize}

\resumeSubheading
{Florida Institute of Technology, Melbourne, FL}
{May 2025 -- Present}
{Graduate Research Assistant}
{Advisor: Dr. Georgios C. Anagnostopoulos}
\begin{itemize}
    \item Built scalable data pipelines to download, organize, and preprocess $\sim$10 years of IMERG precipitation and IR satellite data into analysis-ready geospatial-temporal tensors, reducing manual preparation effort by 70\%.
    \item Designed and trained \textbf{PyTorch} generative models (VAE, Diffusion, Latent Diffusion) for nowcasting/flash-flood prediction, compressing high-dimensional sequences and reducing storage/memory by 35\%.
    \item Implemented end-to-end multimodal forecasting workflows (preprocessing, training, checkpointing, evaluation, experiment logging) and executed multi-GPU \textbf{Slurm} runs to accelerate iteration by 50\%.
    \item Implemented FlowCast radar nowcasting benchmark using Conditional Flow Matching with multi-GPU/multi-node Slurm workflows for scalable model training and evaluation.
\end{itemize}

\resumeSubheading
{Matrixhive Technologies Pvt. Ltd., Ahmedabad, India}
{January 2024 -- April 2024}
{Python Developer Intern}
{}
\begin{itemize}
    \item Developed Django backend modules and RESTful APIs for client-facing workflows; delivered 5+ production features across authentication, dashboards, forms, and admin tools.
    \item Optimized database queries and integrated Redis caching, reducing API latency by 25\% and improving application responsiveness.
\end{itemize}

% ---------- Projects ----------
\section*{Projects}

\resumeProjectHeading
{PKCast: Probabilistic Radar Nowcasting with Conditional Flow Matching}
{\href{https://github.com/prerakpatel51/Pkcast.git}{GitHub Repository}}
{2026}
\begin{itemize}
    \item Developed a \textbf{PyTorch}-based probabilistic radar nowcasting system for the SEVIR VIL dataset to forecast future storm activity from historical radar sequences.
    \item Built a VAE-style autoencoder, generated HDF5 latent datasets, and trained a Conditional Flow Matching model with a Cuboid Transformer U-Net backbone for latent-space forecasting.
    \item Implemented distributed GPU training, MLflow tracking, streaming metrics, Cartopy visualizations, and \textbf{Slurm} scripts for scalable cluster training.
\end{itemize}

\resumeProjectHeading
{3D Variational Autoencoders for Satellite IMERG and IR Data}
{\href{https://github.com/prerakpatel51/icp_neuralnetworks_project}{GitHub Repository}}
{2025}
\begin{itemize}
    \item Designed and trained 3D $\beta$-VAE models on large-scale NASA IMERG precipitation and IR satellite data for spatiotemporal representation learning.
    \item Processed over 1TB of geospatial data and optimized multi-GPU HPC workflows, reducing memory usage by 35\% and accelerating training by 50\%.
\end{itemize}

\resumeProjectHeading
{Satellite Image Segmentation for Environmental Monitoring}
{\href{https://github.com/prerakpatel51/satellite-image-segmentor}{GitHub Repository} | \href{https://huggingface.co/spaces/Prerak51/satellite_image_segmentor}{Live Demo}}
{2024}
\begin{itemize}
    \item Built and deployed a U-Net-based satellite image segmentation system; achieved 86\% segmentation accuracy and deployed an interactive Gradio web interface.
\end{itemize}

% ---------- Education ----------
\section*{Education}

\textbf{Master of Science in Computer Engineering} \hfill January 2025 -- December 2026\\
Florida Institute of Technology\\
Melbourne, FL

\vspace{3pt}

\textbf{Bachelor of Engineering in Computer Engineering} \hfill May 2020 -- May 2024\\
Gandhinagar Institute of Technology\\
Gandhinagar, Gujarat, India

\end{document}